\section*{Problem 2}

\partlabel{(a)}
Steady, constant-density flow; all fluxes and forces are per unit span.
$\mathbf n$ points outward.
% Unit 2 L5 notes, PDF p.5 (printed p.8), Eq. (36).
\begin{align*}
0&=\oint_{CS}\rho\mathbf u\cdot\mathbf n\,dA
 =-\rho UH+\rho\int_0^H u(y)\,dy+\rho v_HL,\\
\int_0^H u(y)\,dy
&=\int_0^\delta\frac{Uy}{\delta}\,dy+\int_\delta^H U\,dy
 =\frac{U\delta}{2}+U(H-\delta)
 =U\left(H-\frac{\delta}{2}\right),\\
v_HL&=UH-U\left(H-\frac{\delta}{2}\right)=\frac{U\delta}{2},\\
v_H&=\frac{U\delta}{2L}
 =\frac{30(1.20\times10^{-3})}{2(0.50)}
 =\boxed{+0.0360~\mathrm{m/s}}.
\end{align*}
Since $v_H>0$, the top flow is outward, in the positive $y$ direction.

\partlabel{(b)}
Neglect streamwise body force and viscous tractions on the inlet, outlet,
and top. Equal inlet/outlet pressures give zero net pressure force in $x$.
$F'_{w\to f,x}$ denotes the wall force on the fluid per unit span.
% Unit 2 L5 notes, PDF p.5, Eq. (37); stress traction is Eq. (12), PDF p.2.
\begin{align*}
F'_{w\to f,x}
&=\oint_{CS}\rho u(\mathbf u\cdot\mathbf n)\,dA
 =-\rho U^2H+\rho\int_0^H u^2(y)\,dy+\rho Uv_HL,\\
\int_0^H u^2(y)\,dy
&=\frac{U^2}{\delta^2}\int_0^\delta y^2\,dy+U^2(H-\delta)
 =U^2\left(H-\frac{2\delta}{3}\right),\\
F'_{w\to f,x}
&=\rho U^2\left[-H+H-\frac{2\delta}{3}+\frac{\delta}{2}\right]
 =-\frac{\rho U^2\delta}{6}\\
&=-\frac{(1.20)(30)^2(1.20\times10^{-3})}{6}
 =\boxed{-0.216~\mathrm{N/m}}.
\end{align*}
The top carries streamwise momentum $\rho Uv_HL=0.648~\mathrm{N/m}$.
The wall force on the fluid is upstream. Its average shear magnitude is
\[
\boxed{\overline\tau_w=-\frac{F'_{w\to f,x}}{L}
=\frac{0.216}{0.50}=0.432~\mathrm{Pa}}.
\]
The signed average wall traction on the fluid is $-0.432~\mathrm{Pa}$.

\partlabel{(c)}
% Unit 2 L5 notes, PDF pp.6--8: D and W definitions, Eqs. (47), (52)--(54).
Use $(\nabla\mathbf u)_{ij}=\dfrac{\partial u_i}{\partial x_j}$:
rows are velocity components and columns are coordinate directions.
\begin{gather*}
\frac{\partial u}{\partial x}=-\frac{v_H}{\delta}=-30~\mathrm{s^{-1}},\qquad
\frac{\partial u}{\partial y}=\frac{U}{\delta}=25000~\mathrm{s^{-1}},\\
\frac{\partial v}{\partial x}=0,\qquad
\frac{\partial v}{\partial y}=\frac{v_H}{\delta}=30~\mathrm{s^{-1}},\\
\boxed{\nabla\mathbf u=
\begin{pmatrix}-30&25000&0\\0&30&0\\0&0&0\end{pmatrix}~\mathrm{s^{-1}}}.
\end{gather*}
\begin{gather*}
\mathbf D=\frac12\left[\nabla\mathbf u+(\nabla\mathbf u)^T\right],\qquad
\mathbf W=\frac12\left[\nabla\mathbf u-(\nabla\mathbf u)^T\right],\\
\boxed{\mathbf D=
\begin{pmatrix}-30&12500&0\\12500&30&0\\0&0&0\end{pmatrix}~\mathrm{s^{-1}}},\qquad
\boxed{\mathbf W=
\begin{pmatrix}0&12500&0\\-12500&0&0\\0&0&0\end{pmatrix}~\mathrm{s^{-1}}}.
\end{gather*}
\[
\boxed{\nabla\cdot\mathbf u
=\frac{\partial u}{\partial x}+\frac{\partial v}{\partial y}
+\frac{\partial w}{\partial z}=-30+30+0=0}.
\]

\partlabel{(d)}
% Unit 2 L5 notes, PDF p.8 (printed p.11), Eqs. (58)--(59).
\begin{gather*}
\mathbf D'=\mathbf D-\frac13(\nabla\cdot\mathbf u)\mathbf I=\mathbf D,\\
\boldsymbol\tau=2\mu\mathbf D'+\zeta(\nabla\cdot\mathbf u)\mathbf I
=2\mu\mathbf D,\\
\boxed{\boldsymbol\tau=
\begin{pmatrix}
-0.001080&0.4500&0\\
0.4500&+0.001080&0\\
0&0&0
\end{pmatrix}~\mathrm{Pa}}.
\end{gather*}
The nonzero stresses are $\tau_{xx}$, $\tau_{yy}$, and $\tau_{xy}=\tau_{yx}$.
The bulk term vanishes because $\nabla\cdot\mathbf u=0$.

\partlabel{(e)}
% Unit 2 L6 notes, PDF p.4 (printed p.18), Eqs. (104), (107).
\begin{align*}
\Phi&=\boldsymbol\tau:\nabla\mathbf u
 =\sum_{i,j}\tau_{ij}\frac{\partial u_i}{\partial x_j}\\
&=\tau_{xx}\frac{\partial u}{\partial x}
 +\tau_{xy}\frac{\partial u}{\partial y}
 +\tau_{yx}\frac{\partial v}{\partial x}
 +\tau_{yy}\frac{\partial v}{\partial y}\\
&=(-0.001080)(-30)+(0.4500)(25000)
 +(0.4500)(0)+(0.001080)(30)\\
&=\mu\left(\frac{U}{\delta}\right)^2
 +4\mu\left(\frac{v_H}{\delta}\right)^2
 \simeq\boxed{1.125\times10^4~\mathrm{W/m^3}}.
\end{align*}
\[
\boxed{\dot s_{\mathrm{gen}}=\frac{\Phi}{T}
=\frac{11250.0648}{300}
\simeq37.50~\mathrm{W/(m^3\,K)}}.
\]
This is viscous entropy production per unit volume.

\partlabel{(f)}
% Unit 2 L6 notes, PDF p.4, Eq. (108); specialize the stress in part (d).
\begin{align*}
\tau_{w,\mathrm{local}}
&\approx\mu\left(\frac{\partial u}{\partial y}
 +\frac{\partial v}{\partial x}\right)_{x=L,y=0}
 =\frac{\mu U}{\delta}
 =\boxed{0.450~\mathrm{Pa}},\\
\frac{\tau_{w,\mathrm{local}}-\overline\tau_w}{\overline\tau_w}\times100\%
&=\frac{0.450-0.432}{0.432}\times100\%
 =\boxed{4.17\%}.
\end{align*}

